\documentclass[12pt]{article}
\usepackage[margin=1in]{geometry}
\usepackage{setspace}
\usepackage{times}
\doublespacing

\begin{document}

\begin{flushleft}
Samuel Okafor\\
Professor Delgado\\
HIST 210: The Atlantic World\\
12 November 2026
\end{flushleft}

\begin{center}
The Printing Press and the Reformation: Technology as Catalyst
\end{center}

When Martin Luther posted his ninety-five theses in 1517, he could not have anticipated
how quickly his arguments would spread across the German-speaking lands.\footnote{The
traditional account of the theses being nailed to the church door at Wittenberg has been
questioned; see Iserloh~\cite{iserloh}.} Within weeks, printers in Leipzig, Nuremberg and
Basel had produced copies, and within a few years Luther had become the most published
author in Europe~\cite{pettegree}.

Historians have long debated whether the press caused the Reformation or merely
accelerated it. Elizabeth Eisenstein argued that print created a ``communications
revolution'' that made the spread of reform inevitable~\cite{eisenstein}. Critics have
replied that oral preaching, song and visual propaganda were equally important, and that
literacy rates in the 1520s remained low.\footnote{Estimates for urban literacy range
from 10 to 30 per cent; rural rates were far lower.}

This essay argues for a middle position. Print did not create the demand for reform, but
it changed the \emph{speed} and \emph{scale} at which ideas could travel, and it gave
reformers a medium that the Church could not easily control. Three features of early
print culture support this view: the cheapness of pamphlets, the networks of
independent printers, and the vernacular.

First, the pamphlet, or \emph{Flugschrift}, was short and cheap. A typical Luther tract
of eight to sixteen pages could be printed in a day and sold for the price of a chicken.
Second, printing in the Empire was decentralised: no single authority could license or
suppress every press. Third, by writing in German rather than Latin, Luther addressed
artisans and merchants directly, and they in turn read his words aloud to others.

In conclusion, the printing press was neither the sole cause of the Reformation nor a
mere bystander. It was a catalyst that turned a local academic dispute into a
continental movement.

\begin{thebibliography}{9}
\bibitem{eisenstein} Elizabeth L. Eisenstein, \emph{The Printing Press as an Agent of
Change}. Cambridge: Cambridge University Press, 1979.
\bibitem{iserloh} Erwin Iserloh, \emph{The Theses Were Not Posted}. Boston: Beacon
Press, 1968.
\bibitem{pettegree} Andrew Pettegree, \emph{Brand Luther}. New York: Penguin, 2015.
\end{thebibliography}

\end{document}
